\section{Discussion and Limitations}
\label{sec:discussion}

The results support a narrow claim: at one sixteenth of the RGB pixel count, a thermal
stream plus a person-aware spatial gate improves nighttime pedestrian segmentation
without sacrificing overall parsing. They do not support a universal sensor budget.
Low-resolution thermal is a proxy for a native sensor, and the bicubic round trip
removes high-frequency detail without modeling sensor noise, fixed-pattern noise, or
the optical blur of a real LWIR array. A real $120\times160$ capture could be better or
worse than our simulation, and only a hardware study can settle the difference. The
same caveat applies to the corrupted-RGB study: the corruption family is synthetic and
was chosen to match the gate's training augmentation, so the absolute degraded numbers
should be read as a stress test rather than a field measurement.

The second limitation is coverage. MFNet is a single urban benchmark, and its nighttime
subset is defined by capture metadata rather than by an independent night protocol. We
attempted a transfer study on PST900 \cite{shivakumar2019rgbthermal}, whose subterranean
RGB-T pairs include a survivor class. The pilot was negative: at $120\times160$ the
thermal-only expert reached 45.35 overall mIoU and 37.88 survivor IoU, while the RGB
expert reached 78.92 and 73.75. The survivors in PST900 are small, often low-contrast,
and frequently occluded, so one sixteenth of the RGB grid is not enough to carry their
thermal evidence. We therefore did not train a gate on PST900; doing so would have
manufactured a second dataset result from a setting in which the premise of the method
fails. The honest conclusion is that the budget knee is dataset-dependent, and that a
deployment should measure it rather than inherit ours.

The third limitation is compute. The gate uses two frozen B2 experts and costs 2.6
times the RGB-only latency. The sensor-pixel saving is real, but it is not a
compute saving, and a reviewer is right to ask for a single-model variant. A native
low-resolution thermal branch could share the RGB backbone and reduce the thermal
expert's cost; a distillation of the gate into the RGB expert could remove the second
forward pass. Both are engineering changes that preserve the routing idea, and both are
outside the scope of this paper. We report the full cost rather than comparing a
two-model system to a one-model baseline without qualification.

The person constraint is a structural prior, not a test-tuned classifier. It encodes
the safety requirement that thermal may repair pedestrians but may not overwrite RGB's
higher-resolution evidence on cars, curbs, guardrails, or road markings. The cost is
visible under corrupted RGB: the person-or gate reaches 49.49 overall, below the full
gate's 51.10, because it refuses to use thermal on non-person classes that RGB can no
longer parse. A validation-selected constraint, or a gate conditioned on image quality,
would be cleaner. We tried a confidence threshold on RGB, but the RGB expert is
confidently wrong at night: its mean maximum softmax probability is 0.983 on both clean
and corrupted inputs, so a scalar threshold cannot separate the regimes. An
image-quality detector that reads brightness, contrast, and noise is the natural next
step.

Finally, the gate is selected on a 157-image held-out split that is larger than the
usual validation split but still small. The three-seed standard deviations are tight
for overall mIoU (0.08) and wider for nighttime pedestrian IoU (0.39), which is the
expected pattern when the safety class is rare. The gate-validation split is
scene-disjoint at the image level but not at the sequence level; MFNet frames from the
same drive can appear in both fit and gate-validation. This does not touch the official
test, but it may make the selection slightly optimistic. A sequence-disjoint split
would be the stronger protocol.
